
%! TeX program = lualatex
\documentclass[justified, nobib,a4paper]{tufte-handout}

\title{Numerical solution of linear advection equation}

\author[Wojciech Sadowski]{Wojciech Sadowski}

%\date{28 March 2010} % without \date command, current date is supplied

% \geometry{showframe} % display margins for debugging page layout

\usepackage{fontspec}
% \setmainfont[Renderer=Basic, Numbers=OldStyle, Scale = 1.0]{TeX Gyre Pagella}
% \setsansfont[Renderer=Basic, Scale=0.90]{TeX Gyre Heros}
% \setmonofont[Renderer=Basic]{TeX Gyre Cursor}

% Set up the spacing using fontspec features
\renewcommand\allcapsspacing[1]{{\addfontfeature{LetterSpace=15}#1}}
\renewcommand\smallcapsspacing[1]{{\addfontfeature{LetterSpace=10}#1}}

% setting ptions for included pictures
\usepackage{graphicx} % allow embedded images
\graphicspath{{graphics/}} % set of paths to search for images
\setkeys{Gin}{
  width=\linewidth,
  totalheight=\textheight,
  keepaspectratio
}

\hypersetup{colorlinks}

\usepackage[compatibility=false, font=footnotesize]{caption}
% * options are necessary to fix the bug that autoref was causing,
%   figure was constantly changed to section
\DeclareCaptionFont{blue}{\color{blue}}
\captionsetup{labelfont={blue}}

% color and color palettes
\usepackage{xcolor}

% drawings collor palette
\definecolor{p11}{HTML}{33cfff}
\definecolor{p12}{HTML}{ffa733}
\definecolor{p13}{HTML}{ff4133}


% drawing inside tex
\usepackage{tikz} 
\usetikzlibrary{decorations.markings}
\usetikzlibrary{calc}
\usetikzlibrary{shapes}

\usepackage{pgfplots}
\pgfplotsset{compat=newest}
\usepgfplotslibrary{fillbetween}
\usepackage{xcolor}

% math things
\usepackage{
  amsmath,  % extended mathematics
  amsthm,
  siunitx,  % non-stacked fractions and better unit spacing
  bm,       % bold math symbols
  physics2,  % a set of usefull symbols for differentiation 
  fixdif,
  derivative
}

\usephysicsmodule{ab, nabla.legacy}

% table things
\usepackage{
  booktabs, % book-quality tables
  multicol  % multiple column layout facilities
} 


% writing references
\usepackage{cleveref}  % must be loded after amsmath

% quoting text and friends
\usepackage{fancyvrb}         % extended verbatim environments
\fvset{fontsize=\normalsize}  % default font size for fancy-verbatim environments
\usepackage{
  dirtytalk, % provides \say command for easy quiting 
  nth       % provides \nth command for 1-st, 2-nd, etc
}        


% citing sources
\usepackage{natbib}
\setcitestyle{authoryear}
\bibliographystyle{plainnat}

% misc 
\usepackage{lipsum}   % filler text

\tikzset{
  set arrow inside/.code={\pgfqkeys{/tikz/arrow inside}{#1}},
  set arrow inside={end/.initial=>, opt/.initial=},
  /pgf/decoration/Mark/.style={
    mark/.expanded=at position #1 with
    {
      \noexpand\arrow[\pgfkeysvalueof{/tikz/arrow inside/opt}]{\pgfkeysvalueof{/tikz/arrow inside/end}}
    }
  },
  arrow inside/.style 2 args={
    set arrow inside={#1},
    postaction={
      decorate,decoration={
        markings,Mark/.list={#2}
      }
    }
  },
}

\input{../src/commands.tex}
\input{../src/symbols.tex}
\newtheorem{theorem}{Theorem}
\newtheorem{lemma}{Lemma}

\newcolumntype{L}[2]{>{\raggedleft#2}p{#1\textwidth}}


\newcommand{\Q}[1]{\section{Exercise:~{#1}}}
\newcommand{\Sol}{\subsection{Solution}}

\tikzset{
  jumpdot/.style={mark=*,solid},
  excl/.append style={jumpdot,fill=white},
  incl/.append style={jumpdot,fill=black},
}

\begin{document}

\maketitle% this prints the handout title, author, and date

\tableofcontents

\section{Introduction}

We will be working with a following equation
\[
	\pdv{u}{t} + a\pdv{u}{x} = 0,\quad a=1
\]
with two initial conditions:
\begin{enumerate}
	\item a smooth one:
	      \[
		      u^0 = \alpha \ab(-8x^2),
	      \]
	      illustrated in \cref{fig:smooth};
	\item a discontinuous one
	      \[
		      u^0 = \begin{cases}
			      1 \text{ for } x \in [-0.5, 0.5], \\
			      0  \text{ otherwise.}
		      \end{cases}
	      \]
	      illustrated in \cref{fig:dics}.
\end{enumerate}
We will be working on a domain $x\in[-1, 1]$.

\paragraph{Exercise:} Plot the analytical solution.


\begin{marginfigure}[-15cm]
	\begin{tikzpicture}
		\begin{axis}[
				smooth,
				% area style,
				xlabel=\(x\),
				ylabel={$u_0$},
				width=6.3cm,
				height=5cm,
				yticklabel=\empty,
				ytick=\empty,
				xtick={-1, 0, 1},
				xmin=-1.25,
				xmax=1.25,
				ymin=0,
				ymax=1.25,
				axis x line=middle,
				axis y line=middle,
				y axis line style={-stealth, thick},
				x axis line style={-stealth, thick},
			]
			\addplot[very thick,samples=400] {exp(-15*x^2)};
			% \addplot[very thick, red, domain=-2:0] {-x} node [midway, left, yshift=-1em] {$\lambda_1 = -a$};
		\end{axis}
	\end{tikzpicture}
	\caption{Smooth initial condition.}\label{fig:smooth}
\end{marginfigure}
\begin{marginfigure}[-5cm]
	\begin{tikzpicture}
		\begin{axis}[
				smooth,
				% area style,
				xlabel=\(x\),
				ylabel={$u_0$},
				width=6.3cm,
				height=5cm,
				yticklabel=\empty,
				ytick=\empty,
				xtick={-1, 0, 1},
				xmin=-1.25,
				xmax=1.25,
				ymin=0,
				ymax=1.25,
				axis x line=middle,
				axis y line=middle,
				y axis line style={-stealth, thick},
				x axis line style={-stealth, thick},
			]
			\addplot[very thick, smooth,samples=20,domain=-2:-0.5]{0};
			\addplot[very thick, smooth,samples=20,domain=0.5:2]{0};
			\addplot[excl] coordinates {(-0.5, 0)};
			\addplot[excl] coordinates {(0.5, 0)};
			\addplot[dashed] coordinates { (0.5, 0)  (0.5, 1) };
			\addplot[dashed] coordinates { (-0.5, 0)  (-0.5, 1) };
			\addplot[very thick, smooth,samples=20,domain=-0.5:0.5]{1};
			\addplot[incl] coordinates {(-0.5, 1)};
			\addplot[incl] coordinates {(0.5, 1)};
		\end{axis}
	\end{tikzpicture}
	\caption{Discontinuous initial condition.}\label{fig:dics}
\end{marginfigure}

\section{Numerical grid}


\begin{marginfigure}
	\begin{tikzpicture}
		\begin{axis}[
				smooth,
				% area style,
				xlabel=\(x\),
				ylabel={$t$},
				width=6.3cm,
				height=5cm,
				ytick={ 0.25, 0.5},
				yticklabels={ $t^{n}$, $t^{n+1}$ },
				xtick={0.3 ,0.5, 0.7 },
				xticklabels={ $x_{i-1}$, $x_{i}$, $x_{i+1}$ },
				grid=both,
				grid style={line width=1pt, draw=gray!75},
				% major grid style={line width=.2pt,draw=gray!50},
				xmin=0,
				xmax=1.25,
				ymin=0,
				ymax=1.25,
				axis x line=middle,
				axis y line=middle,
				y axis line style={-stealth, thick},
				x axis line style={-stealth, thick},
			]
			\addplot[stealth-stealth, thick] coordinates {(0.3, 0.6) (0.5, 0.6)} node [above, midway] {$\Delta x$};
			\addplot[stealth-stealth, thick] coordinates {(0.8, 0.25) (0.8, 0.5)} node [right, midway] {$\Delta t$};
		\end{axis}
	\end{tikzpicture}
	\caption{Numerical grid}\label{fig:grid}
\end{marginfigure}

We will represent our equation on numerical mesh (see
\cref{fig:grid}), consisting of $M$ nodes with
positions labelled as $x_i$, where $1 < i < M$. The nodes are equally spaced by
$\Delta x$. Additionally, our solution will be updated only in certain times,
$t^n$ spaced $\Delta t$. The discrete solution of our continuous function $u$,
at time $t_n$ and position $x_i$, will be denoted as $u_i^n$.

\paragraph{Cyclic boundary conditions}
To make our problem easier we will consider a cyclic grid (whatever \say{goes
	out} on the right comes in from the left). This will make our life easier,
because we do not want to deal with more complex boundary conditions. This will be
equivalent to adding additional $0$ and $N+1$ node to our grid, where we will say
that
\[
	u_0 = u_N,\quad u_{N+1} = u_1.
\]
\section{Gradient representation on the grid}
Our equation consists of gradients, and has to be represented on the numerical grid.
First let's start with the spatial derivative.

\paragraph{Taylor series}
If we want to represent a gradient on a grid we can start with a Taylor series
of a function around a point $a$:
\[
	u \approx \sum\limits_{\alpha=0}^\infty \frac{u^{(n)}(a)(x-a)^n}{n!}
\]
Using the Taylor expansion we can now define a \emph{finite difference} by cutting the series:
\[
	u(x) \approx u(a) + (x-a)u'(a) + \underbrace{\mathcal{O}\ab\{{(x-a)}^2\}}_\text{leading error term}
\]
\[
	u'(a) \approx 	\frac{u(x) - u(a)}{x-a} + \mathcal{O}\ab\{x-a\}
\]
We are looking for a value at $a$, and let's say that we know a value at $x=\xi$, and
$\xi>a$. We can then write a \emph{forwards difference}
\[
	u'(a) \approx 	\frac{u(\xi) - u(a)}{\xi-a} + \mathcal{O}\ab\{\xi-a\}
\]
or if we set $\xi$ and $a$ at nodes that are right next to each other, we can write
\[
	u'_i \approx 	\frac{u_{i+1} - u_i}{\Delta x} + \mathcal{O}\ab(\Delta x).
\]
\paragraph{Upwind}
Analogously, we can write backward difference ($\xi < a$):
\[
	u'_i \approx 	\frac{u_{i} - u_{i-1}}{\Delta x} + \mathcal{O}\ab(\Delta x).
\]
This is an important formula, as it effectively prescribe gradient directionally. We know
that the flow is geeing in the right direction ($a=1$), so we know that the information will
propagate to the right, so we pick backward difference as our gradient representation.

\paragraph{Exercise:} Write and expression for \emph{central difference.} It should be second order accurate and contain $u_{i+1}$ and $u_{i-1}$.

\paragraph{Application to the equation}
We use forward difference for time derivative
\[
	\pdv{u}{t} \approx \frac{u^{n+1} - u^n}{\Delta t},
\]
and we can put all of this into our equation
\[
	\frac{u^{n+1} - u^n}{\Delta t}
	+ a \frac{u_{i} - u_{i-1}}{\Delta x} + \mathcal{O}(\Delta t, \Delta x) = 0,
\]
from which we can get explicit update for our variable
\[
	u^{n+1}_i = u_i^n - \frac{a \Delta t}{\Delta x}(u_i^n - u_{i-1}^n)  + \mathcal{O}(\Delta t, \Delta x) .
\]
The fraction $C = \frac{a \Delta t}{\Delta x}$ is called the CFL number.

\section{Von Neumann stability analysis}


\end{document}
